\SourceLeaf{225}%
der halbe Umfang eines Kreises, dessen Radius $=r$ ist:

\[
=\tfrac{1}{2}\pi k\bigl(e^{\frac{r}{k}}-e^{-\frac{r}{k}}\bigr),
\]

\noindent wobei $k$ eine Konstante ist, von der wir aus Erfahrung wissen, dass sie gegenüber allem, was wir messen können, ungeheuer groß sein muss. In Euklids Geometrie wird sie unendlich.

In der \textls{Bildersprache} des Unendlichen müsste man also sagen, dass sich die Umfänge zweier unendlicher Kreise, deren Radien sich um eine endliche Größe unterscheiden, selbst um eine Größe unterscheiden, die zu ihnen in einem endlichen Verhältnis steht.

Darin liegt aber nichts Widersprüchliches, wenn der endliche Mensch sich nicht anmaßt, etwas Unendliches als etwas Gegebenes und mit seiner gewohnten Anschauung Erfassbares betrachten zu wollen.

Sie sehen, dass hier in der Tat die Frage unmittelbar an die Metaphysik rührt.

\bigskip
\rule{4cm}{0.4pt}
\bigskip

\Needspace{7\baselineskip}
\begin{center}
[\textsc{Schumacher} \emph{an} \textsc{Gauss}.\ \emph{Altona, 19. Juli 1831.}]
\end{center}

\smallskip
\{ Meinen herzlichsten Dank, mein teuerster Freund, für Ihren letzten Brief. Ich kann nicht sagen, dass er mich schon überzeugt hätte. Ich glaube, die unendliche Größe nicht als abgeschlossen gebraucht zu haben. Mir scheint, man kann zeigen, dass mit wachsendem Radius der Abstand zwischen den Eckpunkten des Dreiecks immer mehr verschwindet und sich der Grenze des Zusammenfallens beliebig nähert. Sagt man also der Kürze halber, sie fielen für einen unendlichen Radius wirklich zusammen, so ist dies wie gewöhnlich zu verstehen. Daraus folgt, dass sich, in Bezug auf den Umfang, die von den geraden Linien abgeschnittenen
